\reading{CR-1}{Fundamentals of Credit Risk}{DEFER}{1}

\begin{crkeybox}[Learning objectives for this reading]
\begin{enumerate}[label=\textbf{\alph*.},leftmargin=2em]
\item Define credit risk and explain how it arises using examples.
\item Explain the differences between insolvency, default, and bankruptcy.
\item Identify and describe transactions that generate credit risk.
\item Describe the entities that are exposed to credit risk and explain circumstances under which exposure occurs.
\item Discuss the motivations for managing or taking on credit risk.
\end{enumerate}
\end{crkeybox}

% =====================================================================
\lo{CR-1 a}{Define credit risk and explain how it arises using examples.}

\begin{crdefbox}[Credit risk]
Credit risk refers to the likelihood that a creditor will experience financial
losses if a counterparty fails to fulfill its financial obligations. This
failure can stem from factors like an inability or unwillingness to repay
debts, or delays in honoring obligations.
\end{crdefbox}

Credit risk arises in the following ways.

\begin{enumerate}
\item \term{Inability to repay.} Borrower inability to repay leads to credit
risk. Example: a small business unable to repay a bank loan due to insufficient
revenue.
\item \term{Product obsolescence.} Credit risk arises if products become
obsolete, reducing revenue. Example: a VCR manufacturer unable to repay loans
due to a decline in demand.
\item \term{Unforeseen events.} External factors like economic downturns impact
debtors' ability to repay, posing credit risk. Example: the COVID-19 pandemic
causing businesses to shut down and struggle with loan repayments.
\item \term{Disputes or deliberate actions.} Disputes or debtor actions may lead
to credit risk. Example: a borrower disputes loan terms and refuses repayment,
prompting legal action.
\item \term{Sovereign default.} Governments defaulting on debt obligations pose
credit risk to bondholders. Example: a country defaulting on international debt
or converting foreign currency debt, causing losses for bondholders.
\item \term{Long-term contracts.} Longer contract durations increase credit risk
due to potential adverse events over time.
\end{enumerate}

\begin{crkeybox}[What a lender actually assesses]
When lenders assess credit risk they focus on three aspects:
\begin{enumerate}[label=\alph*)]
\item the amount of potential loss;
\item the probability of borrower default; and
\item the potential recovery of funds, and the timeline for repayment in case of default.
\end{enumerate}
\end{crkeybox}

% =====================================================================
\lo{CR-1 b}{Explain the differences between insolvency, default, and bankruptcy.}

A counterparty's inability to pay its financial obligations can be due to
insolvency, default, or bankruptcy. The three are related but distinct.

\begin{crdefbox}[The three states]
\term{Insolvency} refers to a scenario where a counterparty's liabilities exceed
its assets (i.e., it has negative equity). While insolvency and bankruptcy are
related, insolvent entities are not necessarily bankrupt.

\smallskip
\term{Default} describes a scenario where a counterparty fails to meet its
contractual obligations. A common reason for default is the inability or
unwillingness to pay when an obligation is due. It can be caused by insolvency
but does not necessarily lead to bankruptcy.

\smallskip
\term{Bankruptcy} is a legal procedure where an entity, typically in default,
seeks legal protection through a court. In a bankruptcy process the court
negotiates with the entity's management, creditors, and other stakeholders. The
two forms of bankruptcy are dissolution/liquidation (Chapter 7 in the United
States) and restructuring/reorganization (Chapter 11 in the United States).
\end{crdefbox}

\begin{center}
\begin{tikzpicture}[font=\sffamily\footnotesize]
% three state boxes
\node[draw=FRMNavy, fill=PaleNavy, line width=0.7pt, rounded corners=1pt,
      text width=3.2cm, align=center, minimum height=1.5cm] (ins) at (0,0)
      {\textbf{INSOLVENCY}\\[2pt]\scriptsize a balance-sheet state\\\scriptsize liabilities $>$ assets\\\scriptsize (negative equity)};
\node[draw=FRMGold, fill=PaleGold, line width=0.7pt, rounded corners=1pt,
      text width=3.2cm, align=center, minimum height=1.5cm] (def) at (5.2,0)
      {\textbf{DEFAULT}\\[2pt]\scriptsize an event\\\scriptsize failure to meet a\\\scriptsize contractual obligation};
\node[draw=FRMRed, fill=PaleRed, line width=0.7pt, rounded corners=1pt,
      text width=3.2cm, align=center, minimum height=1.5cm] (bank) at (10.4,0)
      {\textbf{BANKRUPTCY}\\[2pt]\scriptsize a legal procedure\\\scriptsize court protection\\\scriptsize sought by the debtor};
% arrows
\draw[-{Stealth[length=5pt]}, FRMGold, line width=0.9pt] (ins) -- node[above,
  fill=white, inner sep=1.5pt, font=\scriptsize]{can cause} (def);
\draw[-{Stealth[length=5pt]}, FRMRed, line width=0.9pt] (def) -- node[above,
  fill=white, inner sep=1.5pt, font=\scriptsize]{may lead to} (bank);
% two outcomes
\node[draw=FRMRule, fill=white, line width=0.6pt, rounded corners=1pt,
      text width=3.1cm, align=center] (ch7) at (8.5,-2.5)
      {\scriptsize\textbf{Dissolution / liquidation}\\\scriptsize Chapter 7 (US)};
\node[draw=FRMRule, fill=white, line width=0.6pt, rounded corners=1pt,
      text width=3.1cm, align=center] (ch11) at (12.3,-2.5)
      {\scriptsize\textbf{Restructuring / reorganization}\\\scriptsize Chapter 11 (US)};
\draw[-{Stealth[length=4pt]}, FRMGrey] (bank.south) -- ++(0,-0.45) -| (ch7.north);
\draw[-{Stealth[length=4pt]}, FRMGrey] (bank.south) -- ++(0,-0.45) -| (ch11.north);
% the non-implication
\draw[-{Stealth[length=5pt]}, dashed, FRMGrey, line width=0.8pt]
  (ins.south) .. controls (2.8,-2.1) and (7.4,-2.1) .. (bank.south west);
\node[fill=white, inner sep=2pt, font=\scriptsize\itshape, text=FRMGrey]
  at (4.3,-1.7) {insolvent is \textbf{not} necessarily bankrupt};
\end{tikzpicture}
\end{center}
\figcap{Insolvency is a state of the balance sheet, default is an event on a
contract, and bankruptcy is a legal process. The chain runs left to right, but
none of the links is automatic: an entity can be insolvent without ever entering
bankruptcy, and can default without being insolvent.}

% =====================================================================
\lo{CR-1 c}{Identify and describe transactions that generate credit risk.}

\begin{center}
\footnotesize
\begin{tabular}{L{3.6cm} L{5.4cm} L{5.6cm}}
\toprule
\thead{Credit type} & \thead{Losses result from} & \thead{Loss type} \\
\midrule
Loaned money & Nonrepayment; slow repayment; dispute/enforcement &
  Face amount; interest; time value of money; frictional costs \\
Lease obligation & Nonpayment &
  Recovery of asset; remarketing costs; difference in conditions \\
Receivables & Nonpayment of goods delivered or service performed &
  Face amount \\
Prepayment for goods or services & Nondelivery; performance on delivery not as
  contracted; slow delivery; dispute/enforcement &
  Replacement cost; incremental operating cost; time value of money; frictional costs \\
Deposits & Nonrepayment & Face amount; time value of money \\
Claim or contingent claim on asset & Nonrepayment/noncollection; slow
  repayment/slow collection; dispute/enforcement &
  Face amount; time value of money; frictional costs \\
Derivative & Default of third party & Replacement cost (mark-to-market value) \\
\bottomrule
\end{tabular}
\end{center}
\figcap{Types of transactions that create credit risk. Note that the loss type
is not always the face amount: for a derivative it is the replacement cost, and
for a lease it is the shortfall after the asset is recovered and remarketed.}

% =====================================================================
\lo{CR-1 d}{Describe the entities that are exposed to credit risk and explain circumstances under which exposure occurs.}

\subsection*{Banks}
\begin{enumerate}
\item \term{Daily operations.} Banks face significant credit risk from
individual and corporate borrowers and through derivatives activities.
\item \term{Collateralized lending.} Instruments like repurchase agreements
expose banks to counterparty default risk, which is partially mitigated by
collateral. However, market volatility can depreciate collateral values, posing
challenges.
\item \term{Derivatives activities.} Banks encounter counterparty credit risk
through derivatives hedges and portfolios, with notable exposure levels
exemplified by JPMorgan Chase's derivatives receivables exceeding \$700 billion
in 2020.
\end{enumerate}
Despite being sophisticated in credit risk management, banks have seen a decline
in overall risk appetite.

\subsection*{Asset managers}
\begin{enumerate}
\item \term{Investment objectives.} Asset managers cater to varying client risk
objectives, spanning from low-return, low-risk to high-return, high-risk
investments.
\item \term{Client exposure.} Asset managers bear credit risk on behalf of
clients, necessitating thorough risk assessment and oversight.
\item \term{Risk management strategies.} Asset management firms deploy risk
management teams to analyze and mitigate credit risks associated with bond,
equity, and other security investments.
\item \term{Creditworthiness analysis.} Evaluating the creditworthiness of
issuers is integral to risk management, aiding in reducing credit risk exposure.
\end{enumerate}

\subsection*{Hedge funds}
Hedge funds exhibit a higher risk tolerance and pursue aggressive investment
mandates, including risky financial instruments like private equity and debt.
\begin{itemize}
\item \term{Investment strategies.} Strategies involve selling protection
against credit declines or making long and short investments in distressed
securities. Unlike asset managers, hedge fund managers view default as an
investment opportunity.
\item \term{Example.} Hedge funds may short sell securities of distressed
companies or utilize derivatives like credit default swaps (CDSs) for indirect
returns. Notably, Melvin Capital's short sale transactions on GameStop Corp.\ in
2021 led to significant losses when GameStop's share price unexpectedly surged.
\end{itemize}
Significant investments in securities expected to decline entail substantial
risk.

\subsection*{Pension funds}
\begin{itemize}
\item \term{Role and responsibility.} Similar to asset managers, pension funds
invest funds on behalf of pension plan members. Investments in credit-risky
securities can pose significant credit risk for members.
\item \term{Regulatory mandates.} Increasing pension regulation, both for
private and public pension funds, mandates rigorous risk management to protect
plan members.
\end{itemize}

\subsection*{Insurance companies}
\begin{enumerate}
\item \term{Underwriting.} Insurers collect premiums, invest in varied assets,
and pay out claims, balancing income generation with risk minimization.
\item \term{Investments.} They manage client investments separately, shielding
shareholder assets but risking reputation and future business from losses.
\item \term{Reinsurance.} Risks are transferred to reinsurance companies, but
credit risk persists due to claim processing delays and significant claims. Long
time gaps between premium collection and claim settlement pose substantial
credit risk, recorded as reinsurance recoverables on balance sheets.
\end{enumerate}

\subsection*{Corporations}
\begin{enumerate}
\item \term{Account receivables.} Selling goods with deferred payment creates
credit risk as customers may fail to pay. Mitigation strategies include
insurance, factoring, and documentary credit for foreign transactions.
\item \term{Short-term investments and bank deposits.} Investing in short-term
securities and holding deposits in banks entail credit risk due to potential
issuer or bank failures. Diversification across multiple banks helps mitigate
this risk.
\item \term{Derivatives.} Corporations hedge price risks with commodity
derivatives, with over-the-counter contracts posing credit risk in case of
counterparty default. Sophisticated risk management functions analyze and
mitigate such exposures.
\item \term{Vendor financing.} Corporations offering vendor financing face
credit risk from customer default or nonpayment.
\item \term{Supply chain.} Relying on single suppliers or significant shipments
poses credit risk from supplier default or shipment issues. Examples include the
2021 Suez Canal incident causing significant losses in maritime traffic.
\end{enumerate}

\subsection*{Individuals}
People can face credit risk by prepaying for goods or services that might not be
delivered. Bank failures or issues with investment managers can also lead to
losses.

\begin{crkeybox}[The seven exposed entities at a glance]
\footnotesize
\begin{tabular}{L{3.1cm} L{11.6cm}}
\thead{Entity} & \thead{Principal circumstance under which exposure occurs} \\[2pt]
\hline
Banks & Lending to individuals and corporates; repo and other collateralized
lending; derivatives counterparty exposure \\
Asset managers & Credit risk borne \emph{on behalf of clients} across bond,
equity and other holdings \\
Hedge funds & Deliberately taken: selling protection, long/short distressed
positions, CDS \\
Pension funds & Investing plan members' money in credit-risky securities \\
Insurance companies & Underwriting, invested assets, and reinsurance
recoverables over long settlement gaps \\
Corporations & Receivables, deposits and short-term investments, OTC
derivatives, vendor financing, supply chain \\
Individuals & Prepayment for undelivered goods or services; bank or investment
manager failure \\
\end{tabular}
\end{crkeybox}

% =====================================================================
\lo{CR-1 e}{Discuss the motivations for managing or taking on credit risk.}

Credit risk arises out of the consequences of company and management decisions.
Because the company is in control of these decisions, it represents a
\term{controllable risk exposure}.

Successful companies maintain a sufficient equity buffer to absorb some of the
anticipated and unanticipated losses, while considering the company's survival,
profitability, and return on equity:

\begin{enumerate}[label=\Alph*.]
\item \term{Survival.} Managing credit risk ensures that companies avoid large
losses, and therefore, do not face bankruptcy.
\item \term{Profitability.} Managing credit exposure and avoiding losses will
help increase profitability.
\item \term{Return on equity.} Successful companies find the right balance
between debt and equity to maximize their return on equity.
\end{enumerate}
